% GENERATED FILE. Edit canonical lessons, then run npm run generate:books.
% canonical-source-hash: fnv1a64:507e280612ae0103
% canonical-lessons: GU-C263-bhabhi, GU-C263-masi, GU-C263-phoi, GU-C263-sasu, GU-C263-bhatrijo

\chapter{Aunts and in-laws}
\label{ch:gu-aunts-and-in-laws}
\hlchaptermodality{263}

\begin{chapteropening}
\textbf{By the end of this chapter:} \emph{I can ask and answer questions about aunts, in-laws and nephews.}

Complete the last lesson of chapter 263: I can ask and answer questions about aunts, in-laws and nephews.
\end{chapteropening}

\section[\texorpdfstring{\gu{ભાભી} (bhābhī)}{ભાભી (bhābhī)}]{\gu{ભાભી} (bhābhī) — sister-in-law, brother's wife}

\label{lesson:GU-C263-bhabhi}



\begin{quote}
Before the new one: say the Gujarati for picnic, then the Gujarati for to dive.
\end{quote}

\subsection*{You'll want to know: \gu{ભાભી}}
\textbf{\gu{ભાભી}} — \emph{bhābhī} — \textquotedblleft{}sister-in-law, brother's wife\textquotedblright{}.

\begin{quote}
\textbf{\gu{શું} \gu{તમારી} \gu{ભાભી} \gu{અહીં} \gu{રહે} \gu{છે}?} — \emph{śũ tamārī bhābhī ahī̃ rahe che?} — \textquotedblleft{}Does your sister-in-law live here?\textquotedblright{}
\end{quote}

\begin{cousinweb}
The origin of \textbf{\gu{ભાભી}} (\emph{bhābhī}) is not settled. CDIAL lists the older Gujarati \emph{bhojāī}, \textquotedblleft{}brother's wife\textquotedblright{}, under Sanskrit \emph{bhrāturjāyā} (CDIAL 9660), but not \emph{bhābhī}.
\end{cousinweb}

\begin{grammarlens}[title={What you've built}]
One more word for everyday life at A2.
\end{grammarlens}

\subsection*{Your turn}
\begin{itemize}
\raggedright
  \item \emph{Say it:} \emph{bhābhī}
  \item \emph{Say it:} \emph{bhābhī}, once more
  \item \emph{Say it:} \emph{śũ tamārī bhābhī ahī̃ rahe che?} — Does your sister-in-law live here?
  \item \emph{Recall:} say the Gujarati for picnic, then the Gujarati for to dive, then say \emph{bhābhī} again
\end{itemize}

\begin{tcolorbox}[breakable,colback=teal!4,colframe=teal!35!black,title={Before you move on}]
What does \gu{ભાભી} mean? (\textquotedblleft{}Sister-in-law, brother's wife\textquotedblright{}.) Say it once more.
\end{tcolorbox}
\section[\texorpdfstring{\gu{માસી} (māsī)}{માસી (māsī)}]{\gu{માસી} (māsī) — maternal aunt}

\label{lesson:GU-C263-masi}



\begin{quote}
Before the new one: say the Gujarati for to dive, then the Gujarati for sister-in-law.
\end{quote}

\subsection*{You'll want to know: \gu{માસી}}
\textbf{\gu{માસી}} — \emph{māsī} — \textquotedblleft{}maternal aunt\textquotedblright{}.

\begin{quote}
\textbf{\gu{મારી} \gu{માસી} \gu{શહેરમાં} \gu{નથી} \gu{રહેતી}.} — \emph{mārī māsī śahermā̃ nathī rahetī} — \textquotedblleft{}My aunt does not live in the city.\textquotedblright{}
\end{quote}

Your mother's sister.

\begin{cousinweb}
\textbf{\gu{માસી}} (\emph{māsī}) is a tadbhava, a word inherited from Sanskrit and reshaped by sound change. It comes from \emph{mātuḥṣvasṛ}, \textquotedblleft{}mother's sister\textquotedblright{} (CDIAL 10001).
\end{cousinweb}

\begin{grammarlens}[title={What you've built}]
One more word for everyday life at A2.
\end{grammarlens}

\subsection*{Your turn}
\begin{itemize}
\raggedright
  \item \emph{Say it:} \emph{māsī}
  \item \emph{Say it:} \emph{māsī}, once more
  \item \emph{Say it:} \emph{mārī māsī śahermā̃ nathī rahetī} — My aunt does not live in the city.
  \item \emph{Recall:} say the Gujarati for to dive, then the Gujarati for sister-in-law, then say \emph{māsī} again
\end{itemize}

\begin{tcolorbox}[breakable,colback=teal!4,colframe=teal!35!black,title={Before you move on}]
What does \gu{માસી} mean? (\textquotedblleft{}Maternal aunt\textquotedblright{}.) Say it once more.
\end{tcolorbox}
\section[\texorpdfstring{\gu{ફોઈ} (phoī)}{ફોઈ (phoī)}]{\gu{ફોઈ} (phoī) — paternal aunt}

\label{lesson:GU-C263-phoi}



\begin{quote}
Before the new one: say the Gujarati for sister-in-law, then the Gujarati for maternal aunt.
\end{quote}

\subsection*{You'll want to know: \gu{ફોઈ}}
\textbf{\gu{ફોઈ}} — \emph{phoī} — \textquotedblleft{}paternal aunt\textquotedblright{}.

\begin{quote}
\textbf{\gu{શું} \gu{ફોઈ} \gu{આજે} \gu{આવશે}?} — \emph{śũ phoī āje āvśe?} — \textquotedblleft{}Will Aunt come today?\textquotedblright{}
\end{quote}

Your father's sister.

\begin{cousinweb}
\textbf{\gu{ફોઈ}} (\emph{phoī}) is a tadbhava. It comes from \emph{*phupphu-}, \emph{*phupphī-}, \textquotedblleft{}father's sister\textquotedblright{}, reconstructed forms (CDIAL 9089).
\end{cousinweb}

\begin{grammarlens}[title={What you've built}]
One more word for everyday life at A2.
\end{grammarlens}

\subsection*{Your turn}
\begin{itemize}
\raggedright
  \item \emph{Say it:} \emph{phoī}
  \item \emph{Say it:} \emph{phoī}, once more
  \item \emph{Say it:} \emph{śũ phoī āje āvśe?} — Will Aunt come today?
  \item \emph{Recall:} say the Gujarati for sister-in-law, then the Gujarati for maternal aunt, then say \emph{phoī} again
\end{itemize}

\begin{tcolorbox}[breakable,colback=teal!4,colframe=teal!35!black,title={Before you move on}]
What does \gu{ફોઈ} mean? (\textquotedblleft{}Paternal aunt\textquotedblright{}.) Say it once more.
\end{tcolorbox}
\section[\texorpdfstring{\gu{સાસુ} (sāsu)}{સાસુ (sāsu)}]{\gu{સાસુ} (sāsu) — mother-in-law}

\label{lesson:GU-C263-sasu}



\begin{quote}
Before the new one: say the Gujarati for maternal aunt, then the Gujarati for paternal aunt.
\end{quote}

\subsection*{You'll want to know: \gu{સાસુ}}
\textbf{\gu{સાસુ}} — \emph{sāsu} — \textquotedblleft{}mother-in-law\textquotedblright{}.

\begin{quote}
\textbf{\gu{તેની} \gu{સાસુ} \gu{ગુસ્સો} \gu{નથી} \gu{કરતી}.} — \emph{tenī sāsu gusso nathī kartī} — \textquotedblleft{}Her mother-in-law does not get angry.\textquotedblright{}
\end{quote}

\begin{cousinweb}
\textbf{\gu{સાસુ}} (\emph{sāsu}) is a tadbhava. It comes from Sanskrit \emph{śvaśrū́}, \textquotedblleft{}husband's or wife's mother\textquotedblright{} (CDIAL 12759).
\end{cousinweb}

\begin{grammarlens}[title={What you've built}]
One more word for everyday life at A2.
\end{grammarlens}

\subsection*{Your turn}
\begin{itemize}
\raggedright
  \item \emph{Say it:} \emph{sāsu}
  \item \emph{Say it:} \emph{sāsu}, once more
  \item \emph{Say it:} \emph{tenī sāsu gusso nathī kartī} — Her mother-in-law does not get angry.
  \item \emph{Recall:} say the Gujarati for maternal aunt, then the Gujarati for paternal aunt, then say \emph{sāsu} again
\end{itemize}

\begin{tcolorbox}[breakable,colback=teal!4,colframe=teal!35!black,title={Before you move on}]
What does \gu{સાસુ} mean? (\textquotedblleft{}Mother-in-law\textquotedblright{}.) Say it once more.
\end{tcolorbox}
\section[\texorpdfstring{\gu{ભત્રીજો} (bhatrījo)}{ભત્રીજો (bhatrījo)}]{\gu{ભત્રીજો} (bhatrījo) — nephew, brother's son}

\label{lesson:GU-C263-bhatrijo}



\begin{quote}
Before the new one: say the Gujarati for paternal aunt, then the Gujarati for mother-in-law.
\end{quote}

\subsection*{You'll want to know: \gu{ભત્રીજો}}
\textbf{\gu{ભત્રીજો}} — \emph{bhatrījo} — \textquotedblleft{}nephew, brother's son\textquotedblright{}.

\begin{quote}
\textbf{\gu{શું} \gu{તમારો} \gu{ભત્રીજો} \gu{શાળાએ} \gu{જાય} \gu{છે}?} — \emph{śũ tamāro bhatrījo śāḷāe jāy che?} — \textquotedblleft{}Does your nephew go to school?\textquotedblright{}
\end{quote}

\begin{cousinweb}
\textbf{\gu{ભત્રીજો}} (\emph{bhatrījo}) is a tadbhava. It comes from Sanskrit \emph{bhrātr\'{\={\i}}ya}, \textquotedblleft{}brother's son\textquotedblright{}, built on \emph{bhr\'{\={a}}tṛ}, \textquotedblleft{}brother\textquotedblright{} (CDIAL 9672).
\end{cousinweb}

\begin{grammarlens}[title={What you've built}]
One more word for everyday life at A2.
\end{grammarlens}

\subsection*{Your turn}
\begin{itemize}
\raggedright
  \item \emph{Say it:} \emph{bhatrījo}
  \item \emph{Say it:} \emph{bhatrījo}, once more
  \item \emph{Say it:} \emph{śũ tamāro bhatrījo śāḷāe jāy che?} — Does your nephew go to school?
  \item \emph{Recall:} say the Gujarati for paternal aunt, then the Gujarati for mother-in-law, then say \emph{bhatrījo} again
\end{itemize}

\begin{tcolorbox}[breakable,colback=teal!4,colframe=teal!35!black,title={Before you move on}]
What does \gu{ભત્રીજો} mean? (\textquotedblleft{}Nephew, brother's son\textquotedblright{}.) Say it once more.
\end{tcolorbox}
